\documentclass{article}
\usepackage{amsmath, amssymb, amsthm}
\usepackage{hyperref}
\usepackage{geometry}
\geometry{margin=1in}

\title{Erd\H{o}s Problem \#241}
\date{Last updated: 2025-08-31}
\author{Source: erdosproblems.com}

\begin{document}
\maketitle

\noindent\textbf{Status:} OPEN \quad \textbf{Prize: \$100}

\noindent\textbf{Tags:} additive combinatorics, sidon sets

\medskip
\noindent\textbf{Problem Statement:}

Let $f(N)$ be the maximum size of $A\subseteq \{1,\ldots,N\}$ such that the sums $a+b+c$ with $a,b,c\in A$ are all distinct (aside from the trivial coincidences). Is it true that\[ f(N)\sim N^{1/3}?\]

\bigskip
\noindent\textbf{Remarks:}

Originally asked to Erd\H{o}s by Bose. Bose and Chowla \cite{BoCh62} provided a construction proving one half of this, namely\[(1+o(1))N^{1/3}\leq f(N).\]The best upper bound known to date is due to Green \cite{Gr01},\[f(N) \leq ((7/2)^{1/3}+o(1))N^{1/3}\](note that $(7/2)^{1/3}\approx 1.519$). 

More generally, Bose and Chowla conjectured that the maximum size of $A\subseteq \{1,\ldots,N\}$ with all $r$-fold sums distinct (aside from the trivial coincidences) then\[\lvert A\rvert \sim N^{1/r}.\]This is known only for $r=2$ (see [30]).

This is discussed in problem C11 of Guy's collection \cite{Gu04}.

\bigskip
\noindent\textbf{References:}

[BoCh62] Bose, R. C. and Chowla, S., Theorems in the additive theory of numbers. Comment. Math. Helv. (1962/63), 141-147.
 
 [Gr01] Green, Ben, The number of squares and {$B_h[g]$} sets. Acta Arith. (2001), 365-390.
 
 [Gu04] Guy, Richard K., Unsolved problems in number theory. (2004), xviii+437.

\bigskip
\noindent\small{Source: \url{https://www.erdosproblems.com/241}}
\end{document}
